% Generated draft: requires mathematical and source review.
\documentclass[11pt]{article}
\usepackage[margin=25mm]{geometry}
\usepackage{fontspec}
\setmainfont{Latin Modern Roman}
\usepackage{amsmath,amssymb,mathtools,unicode-math}
\setmathfont{Latin Modern Math}
\usepackage{enumitem,needspace,xurl,hyperref}
\setlength{\parindent}{0pt}
\setlength{\parskip}{5pt}
\setlength{\emergencystretch}{3em}
\newcommand{\N}{\mathbb{N}}
\newcommand{\Z}{\mathbb{Z}}
\newcommand{\Q}{\mathbb{Q}}
\newcommand{\R}{\mathbb{R}}
\newcommand{\C}{\mathbb{C}}
\newcommand{\F}{\mathbb{F}}
\newcommand{\A}{\mathbb{A}}
\newcommand{\PP}{\mathbb{P}}
\newcommand{\E}{\mathbb{E}}
\newcommand{\eps}{\varepsilon}
\newcommand{\dd}{\,\mathrm{d}}
\begin{document}
\section*{UnsolvedMath ID: 30001148}
\textbf{Source problem code:} catalogue entry
\subsection{Definitions and mathematical statement}
Let \(K\) be a field complete with respect to a nontrivial discrete valuation, with valuation ring \(\mathcal O_K\), maximal ideal \(\mathfrak m_K\), and residue field \(k=\mathcal O_K/\mathfrak m_K\). Let \(F=K(C)\) be the function field of an integral curve \(C\) over \(K\). No restriction on the characteristic of \(K\) or \(k\) is stated.

Write \(\Omega_F\) for the collection of nontrivial rank-one discrete valuations of \(F\), taken up to equivalence. For \(v\in\Omega_F\), let \(F_v\) denote the completion of \(F\) at \(v\). This interprets ``every discrete valuation'' literally; the supplied record specifies no restriction to valuations arising from a particular model or having a prescribed restriction to \(K\).

Let \(G\) be a connected linear algebraic group over \(F\), and let \(Y\) be an \(F\)-variety equipped with a \(G\)-action making it a homogeneous space: after extension to an algebraic closure \(\overline F\), the action is transitive. For a field extension \(L/F\), write \(Y(L)\) for the set of \(L\)-rational points of \(Y\).

\textbf{Question.} For every such \(K,F,G,Y\), does the implication
\[
\left(\forall v\in\Omega_F,\quad Y(F_v)\ne\varnothing\right)
\quad\Longrightarrow\quad
Y(F)\ne\varnothing
\]
hold?

\textbf{Unresolved qualifications.} The supplied statement imposes no rationality hypothesis on \(G\) or \(Y\), and no additional hypothesis on geometric stabilizers. Its background nevertheless refers to an established result under ``the rationality hypothesis,'' without identifying its subject, its precise meaning, or any accompanying assumptions. That background claim therefore cannot be incorporated as a precise theorem from this record alone. The record also does not specify whether ``curve over \(K\)'' is intended to require geometric integrality or any condition on the constant field of \(F\); the integral-curve formulation above makes the function field meaningful while leaving those additional conditions unresolved.
\subsection{Short English statement}
The question asks whether a homogeneous space of a connected linear algebraic group over the function field of a curve over a complete discretely valued field has a rational point whenever it has a point over every discrete-valuation completion. The supplied record labels this question open, but its intended rationality, stabilizer, and curve qualifications require review.
\subsection{Sources}
\begin{itemize}
\item Local registry: \texttt{data/unsolvedmath/problems.json}, permanent UnsolvedMath ID \(30001148\); supplied source code \texttt{catalogue entry}. The supplied record belongs to the Oberwolfach Reports Open Problems collection and gives proposed year \(2009\); no individual proposer is identified.
\item Supplied primary report citation: \emph{The Arithmetic of Fields} (2009), \emph{Oberwolfach Reports}, pp.~285--354, DOI \texttt{10.4171/owr/2009/05}. \url{https://doi.org/10.4171/owr/2009/05}. The record describes this collection as drawn from mathematical workshop reports and their open-problem discussions.
\item Supplied literature reference: ``Linh, arithmetic of homogeneous spaces.'' \url{https://doi.org/10.1112/jlms.12842}. No fuller bibliographic details or precise supporting theorem are supplied.
\item Literature assessment dated exactly \texttt{2026-08-21}: the supplied curation labels the problem open and reports that patching proves the stated local--global principle under a rationality hypothesis, that no general theorem removing that hypothesis was found, and that later work treats further stabilizer classes rather than the full question. This is supplied dated curation, not an independently verified literature assessment; the status and cited claims require independent verification.
\item The registry labels the statement \texttt{corrected\_verified} and describes its processing as meaning-preserving semantic, typographic, and LaTeX cleanup without unsupported source-dependent additions. These are supplied editorial labels, not independent verification of the mathematical formulation.
\end{itemize}
\end{document}
